\clearpage
\section{Análisis de impacto}\label{sec:impacto}
A corto plazo, el impacto esperado se concentra en disponer de una base técnica organizada y una herramienta demostrable que acerque al usuario previsto a escenarios de comportamiento del fuego. La contribución inmediata será facilitar la lectura conjunta de la situación consultada, su evolución estimada y las condiciones bajo las cuales se obtuvo el resultado, en lugar de entregar únicamente archivos de simulación. La evidencia de este aporte deberá provenir de las comprobaciones y tareas de evaluación descritas en la sección~\ref{sec:transicion}, por lo que no se presume una mejora operativa antes de realizarlas. Para el equipo académico, el desarrollo permitirá consolidar una forma reproducible de preparar datos, identificar restricciones del modelo y conectar productos científicos con requisitos de una aplicación, lo cual constituye un resultado de formación e integración tecnológica dentro del alcance definido.

A mediano plazo, la disponibilidad de escenarios consultables podría servir como apoyo a la capacitación y a la discusión de situaciones de propagación con personas vinculadas a la atención de incendios. Su utilidad social esperada radica en favorecer una comprensión compartida de la exposición y de los límites de la estimación, especialmente cuando las decisiones afectan comunidades próximas al área estudiada. Esta contribución dependerá de que los participantes comprendan la representación y puedan reconocer situaciones que la biblioteca no cubre, de modo que el producto deberá comunicar las restricciones junto con sus resultados. No se compromete una adopción institucional, una sustitución de procedimientos de respuesta ni una disminución cuantificada de tiempos, daños o pérdidas, ya que ninguna de esas consecuencias ha sido evaluada por el proyecto.

A largo plazo, el trabajo podría aportar conocimiento reutilizable sobre las condiciones necesarias para convertir simulaciones físicas en apoyo informativo para la gestión de incendios. En particular, el registro de datos faltantes, resultados desfavorables y límites de aproximación puede orientar a quienes consulten el proyecto y evitar que una representación visual convincente se interprete como certeza científica. Este aporte se relaciona temáticamente con la protección de ecosistemas terrestres y la acción frente a condiciones climáticas consideradas en el póster~\cite{poster2026}, aunque no equivale a demostrar una contribución medible a indicadores de desarrollo sostenible. La proyección se refiere a la utilidad del conocimiento y del producto entregado, sin ampliar el compromiso territorial ni establecer servicios permanentes que excedan la propuesta.
